\section{Заключение}
Контролируемое сравнение восьми способов обучения маршрутизатора
на BANKING77 и CLINC150 установило конкретные режимы сравнительной
деградации и выигрыша. При harm-отборе без коррекции изменение
разности Direct--Weak-error равно \BankShiftEffect{} п.п.
в примере BANKING77 и \ClincShiftEffect{} п.п. в примере
CLINC150, то есть Direct теряет относительное качество.
Во втором датасете найдены также различия внутри
gain-семейства: Hurdle превосходит Direct в harm8 raw на
\HurdleSimpleEffect{} п.п., а взаимодействие метода и состава
составляет \HurdleShiftEffect{} п.п. Все восемь обнаруженных
при accuracy-настройке gain-only взаимодействий устойчивы
к численному уточнению bootstrap.

Коррекция известных вероятностей отбора повышает accuracy Direct
на \CorrectionEffect{} п.п. относительно raw в рассмотренной
CLINC-паре. Дополнительный multiclass-анализ фиксированных решений
устанавливает преимущество Ternary над Hurdle в одном natural-режиме:
\NaturalAccuracyEffect{} п.п. accuracy и \NaturalFoneEffect{} п.п.
macro-F1. Эти результаты дают положительное содержание выбору
способа обучения: важны не только средние на естественном составе,
но и чувствительность реализованного оценивателя к формированию
разметки и корректность сравнительных контролей.

Теоретическое условие смены top-$k$ связывает возможный отказ
с условными harm-вероятностями и зазорами около границы квоты.
Вместе с эмпирическими эффектами оно задаёт направление дальнейшего
исследования: сопоставить одинаковые функциональные классы
и оценить заранее заданное правило выбора на новых задачах.
